\documentclass[svgnames]{article}     % use "amsart" instead of "article" for AMSLaTeX format
%\geometry{landscape}                 % Activate for rotated page geometry

%\usepackage[parfill]{parskip}        % Activate to begin paragraphs with an empty line rather than an indent

\usepackage{graphicx}                 % Use pdf, png, jpg, or eps§ with pdflatex; use eps in DVI mode

%maths                                % TeX will automatically convert eps --> pdf in pdflatex
\usepackage{amssymb}
\usepackage{amsmath}
\usepackage{esint}
\usepackage[top = 0.3in, bottom = 0.3in]{geometry}

% Inverting Color of PDF
%\usepackage{xcolor}
%\pagecolor[rgb]{0.19,0.19,0.19}
%\color[rgb]{0.77,0.77,0.77}

%noindent
\setlength\parindent{0pt}

%pgfplots
\usepackage{pgfplots}

%images
\graphicspath{{/Users/devaldeliwala/research/AQT/review}}                   % Activate to set a image directory

%tikz
\usepackage{pgfplots}
\pgfplotsset{compat=1.15}
\usepackage{comment}
\usetikzlibrary{arrows}
\usepackage[most]{tcolorbox}

%Figures
\usepackage{float}
\usepackage{caption}
\usepackage{lipsum}


\title{Josephson Physics and Interaction Hamiltonian Engineering}
\author{Deval Deliwala}
%\date{}                              % Activate to display a given date or no date

\begin{document}
\maketitle
%\section{}
%\subsection{}
%\tableofcontents                     % Activate to display a table of contents

\section{Josephson Physics}

Here I focus on the physics of the Josephson junction because, being nonlinear,
it is the fundamental circuit element that gives rise to usable qubit states. 
The nonlinearity of the Josephson
junction breaks the degeneracy of the energy level spacings, allowing dynamics
of the system to be restricted to only two qubit states which make up the
\textit{computational basis}. 

A Josephson tunnel is formed by separating two superconducting electrodes with
an insulator thin enough so that electrons can quantum-mechanically
\textit{tunnel} through the barrier, illustrated in Figure \ref{joseph} below. 


\begin{figure}[H]
  \centering
    \includegraphics[width = 6cm]{screenshot 7.png}
    \caption{Schematic diagram of a Josephson junction connected to a bias
    voltage $V$. The Josephson current is given by $I_J = I_0 \sin \delta$,
  where $\delta = \phi_L - \phi_R$ is the difference in the superconducting
phase across the junction}
  \label{joseph}
\end{figure}


The Josephson effect describes the supercurrent ($I_J$ ) that flows through the
junction according to the classical equations 

\begin{align} \label{}
  I_J &= I_0\sin \delta \\ V &= \frac{\Phi_0}{2\pi} \frac{d \delta}{d t},
\end{align}\vspace{3px}

where $\Phi_0 = h / 2e$ is the superconducting flux quantum, $I_0$ is the
critical-current parameter of the junction, and $\delta = \phi_L - \phi_R$ and
$V$ are respectively the superconducting phase difference and voltage across
the junction. The dynamical behavior of these two equations can be represented
by the following equation: 

\[
\frac{d I_J}{d t} = I_0 \cos \delta \frac{2\pi}{\Phi_0} V.
\] \vspace{3px}

with $dI_J / dt$ proportional to $V$, this equation describes an inductor. By
defining the Josephson inductance $L_J$ according to the conventional
definition $V = L_J dI_J / dt$, we get 

\vspace{5px}
\begin{tcolorbox}
  \vspace{-10px}
  \begin{align} \label{}
    L_J = \frac{\Phi_0}{2\pi I_0 \cos \delta}. 
  \end{align}
\end{tcolorbox}
\vspace{5px}

The  $1 / \cos \delta$ term reveals this inductance is nonlinear, becoming
large as $\delta \rightarrow \pi /2$ and being negative for $\pi / 2 < \delta
< 3 \pi / 2$.\footnote{Martinis, Osborne, \textit{Superconducting Qubits and the Physics of Josephson
Junctions}}

Inductance is energy conserving. The energy stored in the Josephson junction is
given by the following formula, 

\begin{align} \label{jospephpotential}
  U_J &= \int I_J(t')V(t') dt' \nonumber \\ 
      &= \int I_0 \sin \delta \cdot \frac{\Phi_0}{2\pi} \frac{d \delta }{d t'} dt'
      \nonumber \\ 
      &= \frac{I_0 \Phi_0}{2\pi} \int \sin \delta d\delta \nonumber \\ 
      &= -\frac{I_0\Phi_0}{2\pi} \cos \delta
\end{align}\vspace{3px}


A generalized Josephson inductance can be calculated from taking the second
derivative of $U$, 

\[
\frac{1}{L_J} = \left( \frac{2\pi}{\Phi_0} \right) ^2 \frac{\partial^2
U(\delta)}{\partial \delta^2}  
\] \vspace{3px}


The general form of the Hamiltonian for the Josephson effect is 

 \begin{align} \label{}
  H_J = U_J
\end{align}\vspace{3px}

\section{Interaction Hamiltonian Engineering} 

To be able to produce entanglement between quantum systems, we need an
\textit{interaction Hamiltonian} that connects degrees of freedoms from each
individual system. Here, I discuss the physical coupling mechanism and its
qubit representation. 

\subsubsection{Physical coupling: capacitive and inductive} 

The Hamiltonian of two coupled systems takes the generalized form 

\[
H = H_1 + H_2 + H_\text{int} 
\] \vspace{3px}

where $H_1$ and $H_2$ are the individual respective Hamiltonians, and
$H_\text{int} $ is the interaction Hamiltonian, which couples variables of both
systems. 

In superconducting circuits, what actually causes coupling energy is an
electric and/or a magnetic field. To achieve \textit{capacitive} coupling,
a capacitor is placed between the voltage nodes of two participating circuits,
yielding an interaction Hamiltonian $H_int$ of the form 

\[
H_\text{int} = C_g V_1 V_2
\] \vspace{3px}

where $C_g$ is the coupling capacitance and $V_1(V_2)$ is the voltage operator
of the corresponding voltage node being connected. 

\begin{figure}[H]
  \centering
    \includegraphics[width = 5cm]{screenshot 8.png}
    \caption{Direct capacitive coupling}
\end{figure}

In the limit $C_g \ll C_1, C_2$, we yield the net Hamiltonian 


\begin{align} \label{hami}
  H = \sum_{i=1,2}^{} [4E_{C_i} n_i^2 - E_{J_i} \cos (\phi_i) ] + 4e^2
  \frac{C_g}{C_1C_2} n_1n_2, 
\end{align}\vspace{3px}

The expression in the brackets are the two original Hamiltonians of the
individual qubits and we take $V_i = (2e / C_i) n_i$. To implement the coupling
capacitance $C_g$, one only needs to bring the edges of the capacitor pads into
close proximity. 

In the case of \textit{inductive} coupling, a mutual inductance shared by two
loops is the coupling mechanism, yielding an interaction Hamiltonian 

\[
H_\text{int} = M_{12} I_1I_2
\] \vspace{3px}

where $M_{12}$ is the mutual inductance between the loops and $I_1$ and $I_2$
are the current operators for the currents through the inductors. 

\begin{figure}[H]
  \centering
    \includegraphics[width = 5cm]{screenshot 9.png}
    \caption{Direct inductive coupling}
    \label{couple}
\end{figure}


The net Hamiltonian can be expressed as 

\begin{align} \label{}
  H = &\sum_{i=1, 2}^{} \left[ 4E_{C_i} n_i^2 - E_{J_i} \cos (\phi_i)
  + \frac{1}{2} \frac{\Phi_{L_i}^2}{L_i (1 - K^2)} \right]\nonumber \\
      &- M_{12} (1 - K^2) \frac{\Phi_{L_1}}{L_1 (1 - K^2)}
      \frac{\Phi_{L_2}}{L_2 (1 - K^2)}
\end{align}\vspace{3px}

where the first two terms are the energies associated with the Josephson
junctions, the third term captures the inductor energies, and the fourth term
is the mutual coupling energy of the form $MI_1I_2$. $\Phi_{L_1}$ and
$\Phi_{L_2}$ are the magnetic fluxes associated with the currents flowing
through the respective inductors. $K$ is the unitless mutual coupling factor
defined by $M_{12} = K\sqrt{L_1L_2}$. 

As we did in the previous note with the charge degree of freedom, we normalize
the inductor and externally applied magnetic fluxes -- in this case, by the
reduced quantum unit of flux $\Phi_0 / 2\pi$ -- to define phases $\phi \equiv
\Phi / (\Phi_0 / 2 \pi )$ that are on the same footing as junction phases
$\phi_1, \phi_2$. These phases must sum to zero or an integer multiple of
$2\pi$. If $I_1, I_2$ are counterclockwise as shown in Figure \ref{couple}., 

\[
  \phi_i - \phi_{L_i} = \phi_{e_i} + 2\pi n \quad n = 0, \pm 1, \hdots. 
\] \vspace{3px}

Replacing $\phi_{L_1(L_2)}$ with $\phi_{1(2)} - \phi_{1(e_2)}$ yields the
Hamiltonian 

\begin{align} \label{}
  H &= \sum_{i=1,2}^{} \left[ 4E_{C_i} n_i^2 - E_{J_i} \cos (\phi_i)
    + \frac{1}{2} \frac{\left( \frac{\Phi_0}{2\pi} \right) ^2 (\phi_i
  -\phi_{ei} )^2 }{L_i (1 - K^2)} \right] \nonumber \\ 
    &- M_{12} (1 - K^2) \frac{\left( \frac{\Phi_0}{2\pi} \right) (\phi_1
    - \phi_{e_1}) }{L_1 (1 - K^2)} \frac{\left( \frac{\Phi_0}{2\pi} \right)
  (\phi_2 - \phi_{e_2})}{L_2(1-K^2)}. 
\end{align}\vspace{3px}

To realize a mutual inductance, two looped currents are brought in close
proximity, or possibly even overlap, or could even share the same wire or
a Josephson inductor. 


\end{document}

